\section{结论}

本文提出 SRPA-YOLO，将预测头容量表述为任务特定的表征分配问题。64 通道共享表征同时支持分类与 DFL 定位，零输出初始化的 16 通道分类私有表征仅提供加性分类残差。该结构扩展分类可用的特征容量，并使边界框输出在计算图上直接独立于私有表征。精确结构重参数化将训练期分解转换为每尺度一个 80 通道卷积和一个合并输出投影，形成静态单路径部署图。

在去重 DUT Anti-UAV 测试集上，验证集选出的 checkpoint 取得 95.850\% Precision、92.201\% Recall、95.185\% AP@0.50、67.071\% \mapall{} 和 77.329\% AP@0.75，部署参数量为 2.164 M。相对 YOLOv8n，五项指标分别提高 3.088、8.242、5.057、6.895 和 9.750 个百分点，同时参数量减少 28.0\%。相对 YOLOv8n-P2，SRPA-YOLO 的 Recall、AP@0.50、\mapall{} 和 AP@0.75 进一步提高，表明检测头任务分配与高分辨率预测具有互补作用。各 IoU 阈值下的一致精度优势以及匹配随机种子实验中的较低性能波动，进一步体现了 SRPA-YOLO 在定位与分类细化方面的稳定表现。

在规定的 RTX 5060、FP16、batch 1、前向推理加 NMS 协议下，SRPA-YOLO 两轮配对延迟为 9.7435 和 9.8574 ms，YOLOv8n 为 9.8859 和 9.9893 ms。结合 28.0\% 的参数量缩减，这些结果表明 SRPA-YOLO 在所评估 GPU 软件栈上获得了有利的精度--紧凑性工作点。后续研究将扩展至多类别航拍场景、更多数据集、嵌入式 ONNX/TensorRT 部署，并分析尺度条件私有容量在不同目标尺度分布下的静态融合性质。
